\documentclass[11pt]{article}

\usepackage[utf8]{inputenc}
\usepackage[T1]{fontenc}
\usepackage{amsmath,amssymb,amsfonts,bm}
\usepackage{geometry}
\geometry{margin=1in}

\title{Supplementary Material:\\
MacroIRT --- Interval-Averaged Bayesian Inference\\
with State-Dependent Open-Channel Noise}
\author{}
\date{}

\begin{document}
\maketitle

\newcommand{\E}{\mathbb{E}}
\newcommand{\Var}{\operatorname{Var}}
\newcommand{\diag}{\operatorname{diag}}
\newcommand{\Nch}{N_{\mathrm{ch}}}
\newcommand{\gbar}{\overline{\gamma}_0}
\newcommand{\sbar}{\overline{\sigma^2}_0}
\newcommand{\tildeg}{\widetilde{\gamma^{T}\Sigma}}
\newcommand{\tildev}{\widetilde{v^{T}\Sigma}}
\newcommand{\tildegg}{\widetilde{\gamma^{T}\Sigma\gamma}}
\newcommand{\tildevv}{\widetilde{v^{T}\Sigma v}}

% ============================================================
\section{Setup and Notation}
% ============================================================

MacroIRT extends MacroIR by accounting for state-dependent open-channel
noise via a Taylor/Laplace correction in boundary-state space. We
adopt the MacroIR setup (boundary states $(i_0,i_t)$, conditional
moments $\overline{\Gamma}_{i_0\to i_t}$ and
$\overline{V}_{i_0\to i_t}$, transition matrix $P(t)=e^{Qt}$,
prior $(\mu_0,\Sigma_0)$, $\Nch$ channels) and reuse the boundary
contractions
\begin{align}
  (\gbar)_{i_0}
  &= \sum_{i_t} P_{i_0\to i_t}(t)\,\overline{\Gamma}_{i_0\to i_t},
  \label{eq:gbar}\\
  (\sbar)_{i_0}
  &= \sum_{i_t} P_{i_0\to i_t}(t)\,\overline{V}_{i_0\to i_t}.
  \label{eq:sbar}
\end{align}
The instrument noise variance over the interval is
$\epsilon^2_{0\to t}=\epsilon^2/t+\nu^2$.

The unified tilde operator $\widetilde{(\cdot)}$ is the same
lift--modulate--collapse transformation defined in the MacroIR
supplement (\texttt{supplement\_macroir\_tilde.md}). We use without
re-derivation the vector and scalar specialisations
\begin{align}
  \widetilde{u^{T}\Sigma}
  &= \overline{u}_0^{T}(\Sigma_0-\diag\,\mu_0)
   + (\overline{U}\circ P)^{T}\mu_0,\\
  \widetilde{u^{T}\Sigma w}
  &= \overline{u}_0^{T}(\Sigma_0-\diag\,\mu_0)\overline{w}_0
   + \mu_0^{T}\big[(\overline{U}\circ P)\circ(\overline{W}\circ P)\big]\mathbf{1}.
\end{align}

% ============================================================
\section{Predictive Mean and Variance}
% ============================================================

Predictive mean (unchanged from MacroIR):
\begin{equation}
  \overline{y}^{\mathrm{pred}}_{0\to t}
  = \Nch\,\mu_0\cdot\gbar.
  \label{eq:ypred}
\end{equation}

Predictive variance (three terms; the third $\sbar$ contribution
matches MacroIR's intrinsic-channel-variability term):
\begin{equation}
  V
  = \epsilon^2_{0\to t}
  + \Nch\,\tildegg
  + \Nch\sum_{i_0}\mu_{0,i_0}\,(\sbar)_{i_0}.
  \label{eq:V}
\end{equation}
The IRT correction does \emph{not} modify $V$; it enters only the
measurement update via the effective direction $v$ defined below.

% ============================================================
\section{Energy in the Boundary-Lifted Representation}
% ============================================================

The IRT energy is structurally identical to the MRT energy but with
$\gbar\to$ tilde objects in the cross-covariance vector and the
effective scalar. Define
\begin{align}
  \delta(p) &= \overline{y}^{\mathrm{obs}}_{0\to t} - \Nch(p\cdot\gbar),\\
  V(p)      &= \epsilon^2_{0\to t} + \Nch\,\tildegg(p) + \Nch(p\cdot\sbar).
\end{align}
Note $V(p)$ depends on $p$ through both the $\sbar$ term and the
$\tildegg$ tilde scalar (which involves the prior $\Sigma_0$ and is
constant in $p$ as long as the boundary moments are evaluated at the
prior, but becomes $p$-dependent if one re-evaluates the tilde at the
candidate $p$; cf.\ MacroIR supplement for the standard convention).

The negative log-posterior (energy) is
\begin{equation}
  E(p) = \log V(p) + \frac{\delta(p)^2}{V(p)}
  + (p-\mu^{\mathrm{prior}})\Nch(\Sigma^{\mathrm{prop}})^{-1}(p-\mu^{\mathrm{prior}})^{T}.
  \label{eq:E}
\end{equation}

\subsection{Tilted vector}

\begin{equation}
  v(p) = \gbar + \frac{\delta(p)}{V(p)}\,\sbar.
  \label{eq:vdef}
\end{equation}

\subsection{Gradient and Hessian (boundary-lifted)}

The pointwise gradient and Hessian formulas have the same algebraic
form as MRT \eqref{eq:hess-shape-MRT-ref} but with each occurrence of
$\Sigma^{\mathrm{prop}}\gbar$ replaced by the boundary tilde
$\tildeg^{T}$ and the $vv^{T}$ outer-product replaced by
$\tildev^{T}\tildev$ in the rank-1 reduction. Specifically, the rank-1
Hessian is
\begin{equation}
  \mathbf{H}_{\mathrm{IRT}}(p)
  = 2\Nch\bigl[(\Sigma^{\mathrm{prop}})^{-1}
  + (\Nch/V)\,(\tildev)^{T}(\tildev)/(v^{T}\Sigma^{\mathrm{prop}}v)\,\dots\bigr]
  \label{eq:hess-IRT-shape}
\end{equation}
and we work directly with the boundary-projected forms. The rigorous
projection identity is the IR tilde lemma proved in the MacroIR
supplement; we apply it to $v$ and $\sbar$ in addition to $\gbar$.

% ============================================================
\section{Rank-1 Quasi-Laplace Update (Main Implementation)}
% ============================================================

\paragraph{Effective scalar.}
\begin{equation}
  s = \tildevv.
  \label{eq:s}
\end{equation}

\paragraph{Posterior covariance.}
\begin{equation}
  \boxed{
    \Sigma^{\mathrm{post}}
    = \Sigma^{\mathrm{prop}}
    - \frac{\Nch}{V+\Nch s}\,
      (\tildev)^{T}(\tildev).
  }
  \label{eq:Sigma_post}
\end{equation}

\paragraph{Posterior mean (from the prior expansion point $p_0=\mu^{\mathrm{prior}}$).}
\begin{equation}
  \boxed{
    \mu^{\mathrm{post}}
    = \mu^{\mathrm{prior}}
    + \frac{\delta}{2V}\,\Sigma^{\mathrm{post}}\,(\tildeg + \tildev).
  }
  \label{eq:mu_post}
\end{equation}

These are the IRT counterparts of the MRT formulas
\eqref{eq:Sigma_post}--\eqref{eq:mu_post} of the MRT supplement, with
the tilde lift in the cross-covariance vector and the effective
scalar.

\paragraph{Working scheme.}
For ensembles of $\Nch\sim 10^2$--$10^4$ the rank-1 quasi-Laplace
form is operational; the rank-2 form including $\tfrac12\log V$ is
analogous to the MRT rank-2 case.

% ============================================================
\section{Newton-MAP Iteration}
% ============================================================

The iteration is identical in structure to the MRT scheme, using
\eqref{eq:Sigma_post}--\eqref{eq:mu_post} for the inner update and
recomputing $\delta,V,v,s,\tildeg,\tildev$ at each iterate. As in
the MRT scheme, one Newton step from $p^{(0)}=\mu^{\mathrm{prior}}$ is
typically sufficient; the multi-step Newton form is needed only when
$|\delta|/V$ is large or $\Nch$ small.

% ============================================================
\section{Reduction Lemmas}
% ============================================================

\subsection{$\sbar=0$ recovers MacroIR}

Setting $\sbar\equiv 0$:
$v=\gbar$, $\tildev=\tildeg$, $s=\tildegg$, and \eqref{eq:Sigma_post}
becomes
\[
  \Sigma^{\mathrm{post}}_{\mathrm{IR}}
  = \Sigma^{\mathrm{prop}} - \frac{1}{V}\,\tildeg^{T}\tildeg,
\]
\[
  \mu^{\mathrm{post}}_{\mathrm{IR}}
  = \mu^{\mathrm{prior}} + \frac{1}{V}\,\Sigma^{\mathrm{prop}}\tildeg\,\delta
\]
(the rank-1 Sherman--Morrison factor $1/(V+\Nch s)$ collapses to
$1/V$ because the $\Nch s$ term is absorbed by the prior already
encoded in the tilde). This is the standard MacroIR scalar Kalman
update.

\subsection{Tilde-collapse to MacroMRT}

Replace each tilde with its per-$i_0$ scalar collapse:
\[
  \tildeg \to \Sigma^{\mathrm{prop}}\gbar,
  \quad
  \tildev \to \Sigma^{\mathrm{prop}}v,
  \quad
  \tildevv \to v^{T}\Sigma^{\mathrm{prop}}v,
\]
and drop the $\Nch\tildegg$ third predictive variance term. This
substitution yields the MacroMRT formulas of the MRT supplement
exactly. The IRT $\to$ MRT collapse is the algebraic statement of
"drop the boundary-state lift, keep the σ² Taylor correction".

\subsection{Diamond identity (the four-corner relationship)}

\[
  \begin{array}{c}
    \text{MacroIRT} \\
    {\scriptstyle \sbar=0}\downarrow \\
    \text{MacroIR}
  \end{array}
  \quad
  \begin{array}{c}
    \xrightarrow{\;\text{tilde-collapse}\;} \\[1ex]
    \xrightarrow{\;\text{tilde-collapse}\;}
  \end{array}
  \quad
  \begin{array}{c}
    \text{MacroMRT} \\
    \downarrow{\scriptstyle \sbar=0} \\
    \text{MacroMR (Moffatt 2007)}
  \end{array}
\]
Both reductions commute. The four corners are pairwise incomparable
along orthogonal axes:
\begin{itemize}
  \item Horizontal axis (tilde lift): captures the boundary
        cross-correlation between channel state and current; this is
        the structural innovation of the I family relative to the M
        family.
  \item Vertical axis (Taylor σ² correction): captures the influence
        of state-dependent open-channel noise on the posterior mean
        and covariance.
\end{itemize}
Pairwise contrasts on the diamond therefore decompose the IRT-vs-MR
gain into its two contributions.

% ============================================================
\section{Trust-Region Simplex Shrink}
% ============================================================

Identical to MacroIR (cf.\ \texttt{macro\_ir\_variance\_inflation\_correction.tex})
and MacroMRT supplement: a shrink factor $\alpha^\star\in[0,1]$ keeps
$\mu^{\mathrm{post}}$ inside the simplex floor and is applied
uniformly to both mean and covariance updates; $\log L$ uses the
unmodified $V$.

% ============================================================
\section{Log-Likelihood Contribution}
% ============================================================

\begin{equation}
  \log L
  = -\tfrac{1}{2}\bigl[\log(2\pi V) + \delta^2/V\bigr],
\end{equation}
with $V$ from \eqref{eq:V}.

\end{document}
